% insanity.tex — dense showcase for hebrew-lualatex-pdf (parallel to the browser skill demo)
% Build from docs/:  bash ../scripts/build.sh insanity.tex
% ============================================================================
% preamble.tex — tested LuaLaTeX preamble for Hebrew + math + figures PDFs (v3.23)
% Battle-tested: compiled an 80-page Hebrew quantum-computing summary AND a
% 197-page Hebrew computer-architecture anthology (TikZ diagrams, code listings,
% colored boxes) — both with 0 missing characters.
%
% Usage in your document:
%   % ============================================================================
% preamble.tex — tested LuaLaTeX preamble for Hebrew + math + figures PDFs (v3.23)
% Battle-tested: compiled an 80-page Hebrew quantum-computing summary AND a
% 197-page Hebrew computer-architecture anthology (TikZ diagrams, code listings,
% colored boxes) — both with 0 missing characters.
%
% Usage in your document:
%   % ============================================================================
% preamble.tex — tested LuaLaTeX preamble for Hebrew + math + figures PDFs (v3.23)
% Battle-tested: compiled an 80-page Hebrew quantum-computing summary AND a
% 197-page Hebrew computer-architecture anthology (TikZ diagrams, code listings,
% colored boxes) — both with 0 missing characters.
%
% Usage in your document:
%   \input{preamble.tex}     % (or paste these lines at the top)
%   \begin{document}  ...  \end{document}
%
% Compile with:  lualatex -interaction=nonstopmode doc.tex   (twice, for TOC)
% Requires the Hebrew fonts installed (run scripts/setup_fonts.sh first).
%
% SHORT doc (exam + solution): change `report` -> `article`.
% STUDY doc meant to be read/annotated: 12pt + \linespread{1.3} below are
%   deliberate readability choices (NOT padding) — keep or trim to taste.
% ============================================================================
\documentclass[12pt]{report}          % 12pt: readable for a study guide; 11pt is fine too
\usepackage[no-math]{fontspec}        % (1) ★ no-math — THE critical line (see recipe.md)
\usepackage[bidi=basic]{babel}        % (2) RTL handling built into the engine
\babelprovide[import,main]{hebrew}    % (3) Hebrew = main language, default RTL
\babelprovide[import]{english}        % (4) English secondary, for correct LTR runs
% Latin fallback registered BEFORE \babelfont so a lone English word in Hebrew
% prose does not tofu. Fonts resolved by FAMILY NAME (portable; no hard path).
\directlua{luaotfload.add_fallback("hebfb",{"DejaVuSerif:mode=node;"})}
% v3.22: Frank Ruhl Libre ships NO italic face. Instead of silently-upright
% \emph (v<=3.20) or bold-as-emphasis (v3.21), synthesize a slanted face with
% FakeSlant — real, VISIBLE italic emphasis with correct \emph nesting
% semantics (inner \emph toggles back upright). Verified empirically at
% 300dpi: slant clearly visible, 0 errors, 0 missing glyphs.
\babelfont{rm}[RawFeature={fallback=hebfb},
  ItalicFont={Frank Ruhl Libre}, ItalicFeatures={FakeSlant=0.22},
  BoldItalicFont={Frank Ruhl Libre Bold}, BoldItalicFeatures={FakeSlant=0.22}
  ]{Frank Ruhl Libre}
\babelfont{sf}[RawFeature={fallback=hebfb}]{Heebo}
\babelfont{tt}[RawFeature={fallback=hebfb}]{DejaVu Sans Mono}

\usepackage{amsmath,amssymb,mathtools}
\usepackage[margin=2.3cm,headheight=14pt]{geometry}
\linespread{1.3}                      % readable line spacing for study/annotation docs
\usepackage{xcolor}
\usepackage{mdframed}
\usepackage{listings}                 % code listings (assembly / C / pseudo-code)
\usepackage{tikz}
\usetikzlibrary{arrows.meta,positioning,calc,shapes.geometric,fit,backgrounds}
\usepackage{pgfplots}\pgfplotsset{compat=1.18}
\usepackage{quantikz}                 % quantum circuits (remove if unused)
\usepackage{booktabs}
\usepackage{longtable}
\usepackage{array}
\usepackage{float}                    % [H] = place table exactly here (reference doc: no drifting)
\usepackage{enumitem}
\usepackage{needspace}

% ============================================================================
% ★★ v3.4 — THE DIRECTION-ISOLATION TRIO (the heart of mixing Hebrew + Latin)
% Mental model:
%   \h{...}    Hebrew (RTL) island      — Hebrew text inside an LTR context
%                                          (TikZ nodes, the english `tikzpic`).
%   \en{...}   English (LTR) isolate    — Latin text / parentheses inside a
%                                          Hebrew (RTL) context. KEEPS PARENS
%                                          AND ORDER CORRECT.
%   \code{...} \en + monospace          — code inside Hebrew / inside a frametitle.
% CRITICAL: never wrap Latin in \h{} (it is for HEBREW). Inside a TikZ node,
% Latin in \h{} renders REVERSED (Datapath->htapataD, CPU->UPC) and parens
% mirror. Isolate every Latin/paren chunk with \en{} or \code{}. See
% references/figures-boxes-listings.md.
% ---- TEXT INSIDE MATH (v3.8) ----------------------------------------------
% Put ANY text that sits inside math in a \text{}. The ONLY thing that is
% invalid in math is a bare \h/\en (= \foreignlanguage) used DIRECTLY, outside
% \text -> "\rmfamily invalid in math mode". Inside \text:
%   Hebrew label           -> \text{...}            (bare Hebrew renders fine;
%                              \h is OPTIONAL/harmless: \text{\h{...}} also works
%                              and both render identically + scale in scripts)
%   plain Latin word        -> \text{...}            (no marker needed)
%   Latin WITH ( ) or [ ]   -> \text{\en{...}}       (\en forces LTR; a plain
%                              \text mirrors the parens in RTL -> "Spec (TNR(" )
%   f(Hebrew/text arg)      -> \text{f}(\text{...})  -- keep the parens at MATH
%                              level, NOT \text{f(arg)} (parens inside the box mirror)
%   code in math            -> \mathtt{...}  (or \texttt{}, math-guarded below)
% NEVER write \en{...} or \h{...} straight into $...$ / \[..\] (outside \text);
% always via \text{}. (\mbox{...} also works but does not scale in scripts.)
% ============================================================================
\newcommand{\h}[1]{\foreignlanguage{hebrew}{#1}}
\newcommand{\en}[1]{\foreignlanguage{english}{#1}}

% ============================================================================
% ★★ — fixes for defects that are INVISIBLE in the log and in the checks
% and only surface on a VISUAL RENDER of the PDF. (Each verified empirically.)
% ----------------------------------------------------------------------------
% v3.22: \emph is now REAL slanted emphasis (FakeSlant italic face defined on
% the \babelfont{rm} line above), replacing the v3.21 \emph->\textbf hack: it
% preserves \emph's toggle semantics and keeps bold free for its own role.
% Prefer bold-as-emphasis for your doc? one line: \renewcommand{\emph}[1]{\textbf{#1}}
% page number: force the digits LTR so they never reverse (18 -> 81) in the RTL
%   footer. That reversal is a NON-DETERMINISTIC bidi glitch (observed once in
%   ~80 pp); this LuaLaTeX build is not byte-reproducible, so the defensive
%   LTR wrap is the only reliable guard. Latin digits are always LTR anyway.
\renewcommand{\thepage}{\foreignlanguage{english}{\arabic{page}}}
% emergencystretch: Hebrew does not hyphenate, so an unbreakable LTR island
%   (\en{...} / an inline math run) landing at the margin can leave a small
%   (~2-13pt) overfull with no legal breakpoint. This rescues ONLY those
%   inline-prose tie-breaks; it is a ceiling used only on offending lines and
%   does NOT touch display-math / listing / table overflow (those still surface
%   in the log for the content-level overflow pass). Verified: baseline overfull
%   -> 0 with this set, with no change to pagination or underfull.
\setlength{\emergencystretch}{3em}
% em-dash: --- does NOT become — under Frank Ruhl (the font lacks tlig; verified
%   the width is identical with/without Ligatures=TeX, while the mechanism works
%   on Latin Modern). TYPE THE CHARACTER — (U+2014) DIRECTLY; likewise – (U+2013).
%   check_bidi_figures.py flags a --/--- that is adjacent to Hebrew.
% ============================================================================

% ---- STRUCTURAL DIRECTION-SAFETY for \texttt (v3.7) -------------------------
% A parenthesis, bracket, or $ INSIDE \texttt{} mirrors/migrates in RTL prose
% (MIPS 0($s0)->"(s0$)0", $rs->"rs$", arr[i]->"[arr[i") because \texttt sets the
% font but NOT the direction. Redefining \texttt to LTR-isolate its content fixes
% the WHOLE class at once (parens, brackets, $) so plain \texttt{lw $t0,0($s0)}
% is safe — no need to remember \code{}. Guarded by \ifmmode because
% \foreignlanguage is invalid in math (there, plain monospace is already LTR).
% NB: \texttt{Hebrew} stays tofu — the monospace font has no Hebrew anyway; put
% Hebrew in \h{}/normal text, never in \texttt{}. To disable, delete this block.
\let\origtexttt\texttt
\renewcommand{\texttt}[1]{\ifmmode\origtexttt{#1}\else\foreignlanguage{english}{\origtexttt{#1}}\fi}
\newcommand{\code}[1]{\ifmmode\origtexttt{#1}\else\foreignlanguage{english}{\origtexttt{#1}}\fi}

% ---- a thin decorative rule helper ----
\newcommand{\sep}{\par\vspace{2pt}\noindent\textcolor{cRule}{\rule{\linewidth}{0.6pt}}\par\vspace{4pt}}

% ---- colors ----
\definecolor{cBlue}{HTML}{1F4E79}   \definecolor{cBlueBg}{HTML}{EAF1F8}
\definecolor{cGreen}{HTML}{2E6B4F}  \definecolor{cGreenBg}{HTML}{EAF4EE}
\definecolor{cAmber}{HTML}{9C6A00}  \definecolor{cAmberBg}{HTML}{FBF3E2}
\definecolor{cPurple}{HTML}{5B2C6F} \definecolor{cPurpleBg}{HTML}{F3ECF7}
\definecolor{cTeal}{HTML}{0E6E6E}   \definecolor{cTealBg}{HTML}{E6F4F4}
\definecolor{cRed}{HTML}{8E2B2B}    \definecolor{cRedBg}{HTML}{FBEAEA}
\definecolor{cGray}{HTML}{555555}   \definecolor{cRule}{HTML}{1F4E79}

% ---- math operators ----
\DeclareMathOperator{\Tr}{Tr}
\DeclareMathOperator{\rank}{rank}
\DeclareMathOperator{\diag}{diag}
\DeclareMathOperator{\sgn}{sgn}

% ---- bra-ket (quantikz defines \ket/\bra/\braket; add the rest safely) ----
\providecommand{\ket}[1]{\left|#1\right\rangle}
\providecommand{\bra}[1]{\left\langle#1\right|}
\providecommand{\braket}[2]{\left\langle#1\middle|#2\right\rangle}
\newcommand{\dyad}[1]{\left|#1\right\rangle\!\left\langle#1\right|}
% v3.21: general outer product |a><b| and expectation value <A> — agents reach
% for these by reflex (esp. \expval); quantikz defines neither, so their absence
% threw "Undefined control sequence" in a real build. \providecommand = safe.
\providecommand{\ketbra}[2]{\left|#1\right\rangle\!\left\langle#2\right|}
\providecommand{\expval}[1]{\left\langle#1\right\rangle}
\newcommand{\abs}[1]{\left|#1\right|}
\newcommand{\norm}[1]{\left\lVert#1\right\rVert}
\newcommand{\half}{\tfrac12}
\newcommand{\R}{\mathbb{R}}\newcommand{\C}{\mathbb{C}}\newcommand{\Z}{\mathbb{Z}}
% identity operator. NB: \mathbb{1} silently renders a WRONG glyph (msbm has no
% blackboard digits). Use \mathbb{I} = 𝕀 (works with amssymb, no extra package).
% Do NOT "fix" this back to \mathbb{1}.
\newcommand{\id}{\mathbb{I}}

% ============================================================================
% Colored boxes — mdframed (NOT tcolorbox, which breaks under bidi).
% ★★ v3.4: innerleftmargin/innerrightmargin give code listings room so their
% frame/line-numbers do not spill out the box border (see listing styles below).
% ============================================================================
\newmdenv[linecolor=cBlue,backgroundcolor=cBlueBg,linewidth=1.1pt,
  frametitlebackgroundcolor=cBlue,frametitlefontcolor=white,
  innertopmargin=6pt,innerbottommargin=7pt,skipabove=9pt,skipbelow=9pt,
  innerleftmargin=9pt,innerrightmargin=9pt,roundcorner=3pt]{defbox}   % הגדרה
\newmdenv[linecolor=cGreen,backgroundcolor=cGreenBg,linewidth=1.1pt,
  frametitlebackgroundcolor=cGreen,frametitlefontcolor=white,
  innertopmargin=6pt,innerbottommargin=7pt,skipabove=9pt,skipbelow=9pt,
  innerleftmargin=9pt,innerrightmargin=9pt,roundcorner=3pt]{thmbox}   % משפט/טענה
\newmdenv[linecolor=cAmber,backgroundcolor=cAmberBg,linewidth=1.1pt,
  frametitlebackgroundcolor=cAmber,frametitlefontcolor=white,
  innertopmargin=6pt,innerbottommargin=7pt,skipabove=9pt,skipbelow=9pt,
  innerleftmargin=9pt,innerrightmargin=9pt,roundcorner=3pt]{notebox}  % הערה
\newmdenv[linecolor=cPurple,backgroundcolor=cPurpleBg,linewidth=1.1pt,
  frametitlebackgroundcolor=cPurple,frametitlefontcolor=white,
  innertopmargin=6pt,innerbottommargin=7pt,skipabove=9pt,skipbelow=9pt,
  innerleftmargin=9pt,innerrightmargin=9pt,roundcorner=3pt]{exbox}    % דוגמה
\newmdenv[linecolor=cRed,backgroundcolor=cRedBg,linewidth=1.1pt,
  frametitlebackgroundcolor=cRed,frametitlefontcolor=white,
  innertopmargin=6pt,innerbottommargin=7pt,skipabove=9pt,skipbelow=9pt,
  innerleftmargin=9pt,innerrightmargin=9pt,roundcorner=3pt]{warnbox}  % טעות נפוצה
\newmdenv[linecolor=cTeal,backgroundcolor=cTealBg,linewidth=1.1pt,
  frametitlebackgroundcolor=cTeal,frametitlefontcolor=white,
  innertopmargin=6pt,innerbottommargin=7pt,skipabove=9pt,skipbelow=9pt,
  innerleftmargin=9pt,innerrightmargin=9pt,roundcorner=3pt]{keybox}   % רעיון מרכזי

% ============================================================================
% Unified TABLE presentation: subtle uniform slate frame + auto-numbered
% caption rendered INSIDE the frame. Use:
%   \begin{tablebox}\centering\small  <tabular>  \tabcap{<description>}\end{tablebox}
% Replaces the buggy \caption* (which emitted a stray "טבלה N : *" line) and
% gives every table one identical frame + numbered caption "טבלה N.M.".
% Optional intro text may be placed at the TOP, before the tabular, inside the box.
% ============================================================================
\definecolor{cTblLine}{HTML}{B9C4D2}  \definecolor{cTblBg}{HTML}{FAFBFC}
% v3.23: class-aware — under `article` there is no chapter counter (the promised
% report->article switch used to die with "No counter 'chapter' defined").
\makeatletter
\@ifundefined{c@chapter}{%
  \newcounter{tbl}[section]\renewcommand{\thetbl}{\arabic{tbl}}%
}{%
  \newcounter{tbl}[chapter]\renewcommand{\thetbl}{\thechapter.\arabic{tbl}}%
}
\makeatother
\newmdenv[linecolor=cTblLine,backgroundcolor=cTblBg,linewidth=0.9pt,
  innertopmargin=9pt,innerbottommargin=8pt,skipabove=11pt,skipbelow=11pt,
  innerleftmargin=11pt,innerrightmargin=11pt,roundcorner=4pt]{tablebox}
\newcommand{\tabcap}[1]{\refstepcounter{tbl}\par\addvspace{6pt}%
  {\centering\small\textcolor{cGray}{\textbf{טבלה~\thetbl.}}\enspace #1\par}}

% ============================================================================
% LTR islands.  EVERY tikzpicture / quantikz MUST live inside one, or the whole
% figure's coordinates and Latin labels flip. (Hebrew node text still uses \h{}.)
% ============================================================================
\newenvironment{qcirc}{\par\medskip\begin{otherlanguage}{english}\begin{center}}
  {\end{center}\end{otherlanguage}\par\medskip}
\newenvironment{tikzpic}{\par\medskip\begin{otherlanguage}{english}\begin{center}}
  {\end{center}\end{otherlanguage}\par\medskip}

% ============================================================================
% Code listings.  ★★ v3.4 THE BOX-OVERFLOW FIX:
%   resetmargins=true   → margins measured from the ENCLOSING box, not the page
%                         (without it, a listing inside an mdframed box spills
%                          its frame + line-numbers out past the box border).
%   framexleftmargin=0  → the frame (left bar) does not extend left of the code.
%   xleftmargin keeps a small indent so the line numbers have room INSIDE.
% asmblock wraps the listing in an LTR island (code is LTR).
% ============================================================================
\lstdefinestyle{asm}{
  basicstyle=\ttfamily\small,
  backgroundcolor=\color{cBlueBg!55},
  frame=leftline,framerule=2.5pt,rulecolor=\color{cBlue},
  numbers=left,numberstyle=\tiny\color{cGray},numbersep=8pt,
  xleftmargin=2.2em,framexleftmargin=0pt,framexrightmargin=0pt,resetmargins=true,
  breaklines=true,columns=fullflexible,keepspaces=true,
  showstringspaces=false,aboveskip=8pt,belowskip=8pt,
  commentstyle=\color{cGreen},
}
\lstdefinestyle{cstyle}{
  language=C,basicstyle=\ttfamily\small,
  backgroundcolor=\color{cGreenBg!55},
  frame=leftline,framerule=2.5pt,rulecolor=\color{cGreen},
  numbers=left,numberstyle=\tiny\color{cGray},numbersep=8pt,
  xleftmargin=2.2em,framexleftmargin=0pt,framexrightmargin=0pt,resetmargins=true,
  breaklines=true,columns=fullflexible,keepspaces=true,
  showstringspaces=false,keywordstyle=\color{cBlue}\bfseries,
  commentstyle=\color{cGray}\itshape,aboveskip=8pt,belowskip=8pt,
}
\newenvironment{asmblock}{\par\begin{otherlanguage}{english}}{\end{otherlanguage}\par}

\usepackage[hidelinks,unicode]{hyperref}   % LAST, after colors. unicode = Hebrew bookmarks.

\setcounter{tocdepth}{1}
\setcounter{secnumdepth}{2}
     % (or paste these lines at the top)
%   \begin{document}  ...  \end{document}
%
% Compile with:  lualatex -interaction=nonstopmode doc.tex   (twice, for TOC)
% Requires the Hebrew fonts installed (run scripts/setup_fonts.sh first).
%
% SHORT doc (exam + solution): change `report` -> `article`.
% STUDY doc meant to be read/annotated: 12pt + \linespread{1.3} below are
%   deliberate readability choices (NOT padding) — keep or trim to taste.
% ============================================================================
\documentclass[12pt]{report}          % 12pt: readable for a study guide; 11pt is fine too
\usepackage[no-math]{fontspec}        % (1) ★ no-math — THE critical line (see recipe.md)
\usepackage[bidi=basic]{babel}        % (2) RTL handling built into the engine
\babelprovide[import,main]{hebrew}    % (3) Hebrew = main language, default RTL
\babelprovide[import]{english}        % (4) English secondary, for correct LTR runs
% Latin fallback registered BEFORE \babelfont so a lone English word in Hebrew
% prose does not tofu. Fonts resolved by FAMILY NAME (portable; no hard path).
\directlua{luaotfload.add_fallback("hebfb",{"DejaVuSerif:mode=node;"})}
% v3.22: Frank Ruhl Libre ships NO italic face. Instead of silently-upright
% \emph (v<=3.20) or bold-as-emphasis (v3.21), synthesize a slanted face with
% FakeSlant — real, VISIBLE italic emphasis with correct \emph nesting
% semantics (inner \emph toggles back upright). Verified empirically at
% 300dpi: slant clearly visible, 0 errors, 0 missing glyphs.
\babelfont{rm}[RawFeature={fallback=hebfb},
  ItalicFont={Frank Ruhl Libre}, ItalicFeatures={FakeSlant=0.22},
  BoldItalicFont={Frank Ruhl Libre Bold}, BoldItalicFeatures={FakeSlant=0.22}
  ]{Frank Ruhl Libre}
\babelfont{sf}[RawFeature={fallback=hebfb}]{Heebo}
\babelfont{tt}[RawFeature={fallback=hebfb}]{DejaVu Sans Mono}

\usepackage{amsmath,amssymb,mathtools}
\usepackage[margin=2.3cm,headheight=14pt]{geometry}
\linespread{1.3}                      % readable line spacing for study/annotation docs
\usepackage{xcolor}
\usepackage{mdframed}
\usepackage{listings}                 % code listings (assembly / C / pseudo-code)
\usepackage{tikz}
\usetikzlibrary{arrows.meta,positioning,calc,shapes.geometric,fit,backgrounds}
\usepackage{pgfplots}\pgfplotsset{compat=1.18}
\usepackage{quantikz}                 % quantum circuits (remove if unused)
\usepackage{booktabs}
\usepackage{longtable}
\usepackage{array}
\usepackage{float}                    % [H] = place table exactly here (reference doc: no drifting)
\usepackage{enumitem}
\usepackage{needspace}

% ============================================================================
% ★★ v3.4 — THE DIRECTION-ISOLATION TRIO (the heart of mixing Hebrew + Latin)
% Mental model:
%   \h{...}    Hebrew (RTL) island      — Hebrew text inside an LTR context
%                                          (TikZ nodes, the english `tikzpic`).
%   \en{...}   English (LTR) isolate    — Latin text / parentheses inside a
%                                          Hebrew (RTL) context. KEEPS PARENS
%                                          AND ORDER CORRECT.
%   \code{...} \en + monospace          — code inside Hebrew / inside a frametitle.
% CRITICAL: never wrap Latin in \h{} (it is for HEBREW). Inside a TikZ node,
% Latin in \h{} renders REVERSED (Datapath->htapataD, CPU->UPC) and parens
% mirror. Isolate every Latin/paren chunk with \en{} or \code{}. See
% references/figures-boxes-listings.md.
% ---- TEXT INSIDE MATH (v3.8) ----------------------------------------------
% Put ANY text that sits inside math in a \text{}. The ONLY thing that is
% invalid in math is a bare \h/\en (= \foreignlanguage) used DIRECTLY, outside
% \text -> "\rmfamily invalid in math mode". Inside \text:
%   Hebrew label           -> \text{...}            (bare Hebrew renders fine;
%                              \h is OPTIONAL/harmless: \text{\h{...}} also works
%                              and both render identically + scale in scripts)
%   plain Latin word        -> \text{...}            (no marker needed)
%   Latin WITH ( ) or [ ]   -> \text{\en{...}}       (\en forces LTR; a plain
%                              \text mirrors the parens in RTL -> "Spec (TNR(" )
%   f(Hebrew/text arg)      -> \text{f}(\text{...})  -- keep the parens at MATH
%                              level, NOT \text{f(arg)} (parens inside the box mirror)
%   code in math            -> \mathtt{...}  (or \texttt{}, math-guarded below)
% NEVER write \en{...} or \h{...} straight into $...$ / \[..\] (outside \text);
% always via \text{}. (\mbox{...} also works but does not scale in scripts.)
% ============================================================================
\newcommand{\h}[1]{\foreignlanguage{hebrew}{#1}}
\newcommand{\en}[1]{\foreignlanguage{english}{#1}}

% ============================================================================
% ★★ — fixes for defects that are INVISIBLE in the log and in the checks
% and only surface on a VISUAL RENDER of the PDF. (Each verified empirically.)
% ----------------------------------------------------------------------------
% v3.22: \emph is now REAL slanted emphasis (FakeSlant italic face defined on
% the \babelfont{rm} line above), replacing the v3.21 \emph->\textbf hack: it
% preserves \emph's toggle semantics and keeps bold free for its own role.
% Prefer bold-as-emphasis for your doc? one line: \renewcommand{\emph}[1]{\textbf{#1}}
% page number: force the digits LTR so they never reverse (18 -> 81) in the RTL
%   footer. That reversal is a NON-DETERMINISTIC bidi glitch (observed once in
%   ~80 pp); this LuaLaTeX build is not byte-reproducible, so the defensive
%   LTR wrap is the only reliable guard. Latin digits are always LTR anyway.
\renewcommand{\thepage}{\foreignlanguage{english}{\arabic{page}}}
% emergencystretch: Hebrew does not hyphenate, so an unbreakable LTR island
%   (\en{...} / an inline math run) landing at the margin can leave a small
%   (~2-13pt) overfull with no legal breakpoint. This rescues ONLY those
%   inline-prose tie-breaks; it is a ceiling used only on offending lines and
%   does NOT touch display-math / listing / table overflow (those still surface
%   in the log for the content-level overflow pass). Verified: baseline overfull
%   -> 0 with this set, with no change to pagination or underfull.
\setlength{\emergencystretch}{3em}
% em-dash: --- does NOT become — under Frank Ruhl (the font lacks tlig; verified
%   the width is identical with/without Ligatures=TeX, while the mechanism works
%   on Latin Modern). TYPE THE CHARACTER — (U+2014) DIRECTLY; likewise – (U+2013).
%   check_bidi_figures.py flags a --/--- that is adjacent to Hebrew.
% ============================================================================

% ---- STRUCTURAL DIRECTION-SAFETY for \texttt (v3.7) -------------------------
% A parenthesis, bracket, or $ INSIDE \texttt{} mirrors/migrates in RTL prose
% (MIPS 0($s0)->"(s0$)0", $rs->"rs$", arr[i]->"[arr[i") because \texttt sets the
% font but NOT the direction. Redefining \texttt to LTR-isolate its content fixes
% the WHOLE class at once (parens, brackets, $) so plain \texttt{lw $t0,0($s0)}
% is safe — no need to remember \code{}. Guarded by \ifmmode because
% \foreignlanguage is invalid in math (there, plain monospace is already LTR).
% NB: \texttt{Hebrew} stays tofu — the monospace font has no Hebrew anyway; put
% Hebrew in \h{}/normal text, never in \texttt{}. To disable, delete this block.
\let\origtexttt\texttt
\renewcommand{\texttt}[1]{\ifmmode\origtexttt{#1}\else\foreignlanguage{english}{\origtexttt{#1}}\fi}
\newcommand{\code}[1]{\ifmmode\origtexttt{#1}\else\foreignlanguage{english}{\origtexttt{#1}}\fi}

% ---- a thin decorative rule helper ----
\newcommand{\sep}{\par\vspace{2pt}\noindent\textcolor{cRule}{\rule{\linewidth}{0.6pt}}\par\vspace{4pt}}

% ---- colors ----
\definecolor{cBlue}{HTML}{1F4E79}   \definecolor{cBlueBg}{HTML}{EAF1F8}
\definecolor{cGreen}{HTML}{2E6B4F}  \definecolor{cGreenBg}{HTML}{EAF4EE}
\definecolor{cAmber}{HTML}{9C6A00}  \definecolor{cAmberBg}{HTML}{FBF3E2}
\definecolor{cPurple}{HTML}{5B2C6F} \definecolor{cPurpleBg}{HTML}{F3ECF7}
\definecolor{cTeal}{HTML}{0E6E6E}   \definecolor{cTealBg}{HTML}{E6F4F4}
\definecolor{cRed}{HTML}{8E2B2B}    \definecolor{cRedBg}{HTML}{FBEAEA}
\definecolor{cGray}{HTML}{555555}   \definecolor{cRule}{HTML}{1F4E79}

% ---- math operators ----
\DeclareMathOperator{\Tr}{Tr}
\DeclareMathOperator{\rank}{rank}
\DeclareMathOperator{\diag}{diag}
\DeclareMathOperator{\sgn}{sgn}

% ---- bra-ket (quantikz defines \ket/\bra/\braket; add the rest safely) ----
\providecommand{\ket}[1]{\left|#1\right\rangle}
\providecommand{\bra}[1]{\left\langle#1\right|}
\providecommand{\braket}[2]{\left\langle#1\middle|#2\right\rangle}
\newcommand{\dyad}[1]{\left|#1\right\rangle\!\left\langle#1\right|}
% v3.21: general outer product |a><b| and expectation value <A> — agents reach
% for these by reflex (esp. \expval); quantikz defines neither, so their absence
% threw "Undefined control sequence" in a real build. \providecommand = safe.
\providecommand{\ketbra}[2]{\left|#1\right\rangle\!\left\langle#2\right|}
\providecommand{\expval}[1]{\left\langle#1\right\rangle}
\newcommand{\abs}[1]{\left|#1\right|}
\newcommand{\norm}[1]{\left\lVert#1\right\rVert}
\newcommand{\half}{\tfrac12}
\newcommand{\R}{\mathbb{R}}\newcommand{\C}{\mathbb{C}}\newcommand{\Z}{\mathbb{Z}}
% identity operator. NB: \mathbb{1} silently renders a WRONG glyph (msbm has no
% blackboard digits). Use \mathbb{I} = 𝕀 (works with amssymb, no extra package).
% Do NOT "fix" this back to \mathbb{1}.
\newcommand{\id}{\mathbb{I}}

% ============================================================================
% Colored boxes — mdframed (NOT tcolorbox, which breaks under bidi).
% ★★ v3.4: innerleftmargin/innerrightmargin give code listings room so their
% frame/line-numbers do not spill out the box border (see listing styles below).
% ============================================================================
\newmdenv[linecolor=cBlue,backgroundcolor=cBlueBg,linewidth=1.1pt,
  frametitlebackgroundcolor=cBlue,frametitlefontcolor=white,
  innertopmargin=6pt,innerbottommargin=7pt,skipabove=9pt,skipbelow=9pt,
  innerleftmargin=9pt,innerrightmargin=9pt,roundcorner=3pt]{defbox}   % הגדרה
\newmdenv[linecolor=cGreen,backgroundcolor=cGreenBg,linewidth=1.1pt,
  frametitlebackgroundcolor=cGreen,frametitlefontcolor=white,
  innertopmargin=6pt,innerbottommargin=7pt,skipabove=9pt,skipbelow=9pt,
  innerleftmargin=9pt,innerrightmargin=9pt,roundcorner=3pt]{thmbox}   % משפט/טענה
\newmdenv[linecolor=cAmber,backgroundcolor=cAmberBg,linewidth=1.1pt,
  frametitlebackgroundcolor=cAmber,frametitlefontcolor=white,
  innertopmargin=6pt,innerbottommargin=7pt,skipabove=9pt,skipbelow=9pt,
  innerleftmargin=9pt,innerrightmargin=9pt,roundcorner=3pt]{notebox}  % הערה
\newmdenv[linecolor=cPurple,backgroundcolor=cPurpleBg,linewidth=1.1pt,
  frametitlebackgroundcolor=cPurple,frametitlefontcolor=white,
  innertopmargin=6pt,innerbottommargin=7pt,skipabove=9pt,skipbelow=9pt,
  innerleftmargin=9pt,innerrightmargin=9pt,roundcorner=3pt]{exbox}    % דוגמה
\newmdenv[linecolor=cRed,backgroundcolor=cRedBg,linewidth=1.1pt,
  frametitlebackgroundcolor=cRed,frametitlefontcolor=white,
  innertopmargin=6pt,innerbottommargin=7pt,skipabove=9pt,skipbelow=9pt,
  innerleftmargin=9pt,innerrightmargin=9pt,roundcorner=3pt]{warnbox}  % טעות נפוצה
\newmdenv[linecolor=cTeal,backgroundcolor=cTealBg,linewidth=1.1pt,
  frametitlebackgroundcolor=cTeal,frametitlefontcolor=white,
  innertopmargin=6pt,innerbottommargin=7pt,skipabove=9pt,skipbelow=9pt,
  innerleftmargin=9pt,innerrightmargin=9pt,roundcorner=3pt]{keybox}   % רעיון מרכזי

% ============================================================================
% Unified TABLE presentation: subtle uniform slate frame + auto-numbered
% caption rendered INSIDE the frame. Use:
%   \begin{tablebox}\centering\small  <tabular>  \tabcap{<description>}\end{tablebox}
% Replaces the buggy \caption* (which emitted a stray "טבלה N : *" line) and
% gives every table one identical frame + numbered caption "טבלה N.M.".
% Optional intro text may be placed at the TOP, before the tabular, inside the box.
% ============================================================================
\definecolor{cTblLine}{HTML}{B9C4D2}  \definecolor{cTblBg}{HTML}{FAFBFC}
% v3.23: class-aware — under `article` there is no chapter counter (the promised
% report->article switch used to die with "No counter 'chapter' defined").
\makeatletter
\@ifundefined{c@chapter}{%
  \newcounter{tbl}[section]\renewcommand{\thetbl}{\arabic{tbl}}%
}{%
  \newcounter{tbl}[chapter]\renewcommand{\thetbl}{\thechapter.\arabic{tbl}}%
}
\makeatother
\newmdenv[linecolor=cTblLine,backgroundcolor=cTblBg,linewidth=0.9pt,
  innertopmargin=9pt,innerbottommargin=8pt,skipabove=11pt,skipbelow=11pt,
  innerleftmargin=11pt,innerrightmargin=11pt,roundcorner=4pt]{tablebox}
\newcommand{\tabcap}[1]{\refstepcounter{tbl}\par\addvspace{6pt}%
  {\centering\small\textcolor{cGray}{\textbf{טבלה~\thetbl.}}\enspace #1\par}}

% ============================================================================
% LTR islands.  EVERY tikzpicture / quantikz MUST live inside one, or the whole
% figure's coordinates and Latin labels flip. (Hebrew node text still uses \h{}.)
% ============================================================================
\newenvironment{qcirc}{\par\medskip\begin{otherlanguage}{english}\begin{center}}
  {\end{center}\end{otherlanguage}\par\medskip}
\newenvironment{tikzpic}{\par\medskip\begin{otherlanguage}{english}\begin{center}}
  {\end{center}\end{otherlanguage}\par\medskip}

% ============================================================================
% Code listings.  ★★ v3.4 THE BOX-OVERFLOW FIX:
%   resetmargins=true   → margins measured from the ENCLOSING box, not the page
%                         (without it, a listing inside an mdframed box spills
%                          its frame + line-numbers out past the box border).
%   framexleftmargin=0  → the frame (left bar) does not extend left of the code.
%   xleftmargin keeps a small indent so the line numbers have room INSIDE.
% asmblock wraps the listing in an LTR island (code is LTR).
% ============================================================================
\lstdefinestyle{asm}{
  basicstyle=\ttfamily\small,
  backgroundcolor=\color{cBlueBg!55},
  frame=leftline,framerule=2.5pt,rulecolor=\color{cBlue},
  numbers=left,numberstyle=\tiny\color{cGray},numbersep=8pt,
  xleftmargin=2.2em,framexleftmargin=0pt,framexrightmargin=0pt,resetmargins=true,
  breaklines=true,columns=fullflexible,keepspaces=true,
  showstringspaces=false,aboveskip=8pt,belowskip=8pt,
  commentstyle=\color{cGreen},
}
\lstdefinestyle{cstyle}{
  language=C,basicstyle=\ttfamily\small,
  backgroundcolor=\color{cGreenBg!55},
  frame=leftline,framerule=2.5pt,rulecolor=\color{cGreen},
  numbers=left,numberstyle=\tiny\color{cGray},numbersep=8pt,
  xleftmargin=2.2em,framexleftmargin=0pt,framexrightmargin=0pt,resetmargins=true,
  breaklines=true,columns=fullflexible,keepspaces=true,
  showstringspaces=false,keywordstyle=\color{cBlue}\bfseries,
  commentstyle=\color{cGray}\itshape,aboveskip=8pt,belowskip=8pt,
}
\newenvironment{asmblock}{\par\begin{otherlanguage}{english}}{\end{otherlanguage}\par}

\usepackage[hidelinks,unicode]{hyperref}   % LAST, after colors. unicode = Hebrew bookmarks.

\setcounter{tocdepth}{1}
\setcounter{secnumdepth}{2}
     % (or paste these lines at the top)
%   \begin{document}  ...  \end{document}
%
% Compile with:  lualatex -interaction=nonstopmode doc.tex   (twice, for TOC)
% Requires the Hebrew fonts installed (run scripts/setup_fonts.sh first).
%
% SHORT doc (exam + solution): change `report` -> `article`.
% STUDY doc meant to be read/annotated: 12pt + \linespread{1.3} below are
%   deliberate readability choices (NOT padding) — keep or trim to taste.
% ============================================================================
\documentclass[12pt]{report}          % 12pt: readable for a study guide; 11pt is fine too
\usepackage[no-math]{fontspec}        % (1) ★ no-math — THE critical line (see recipe.md)
\usepackage[bidi=basic]{babel}        % (2) RTL handling built into the engine
\babelprovide[import,main]{hebrew}    % (3) Hebrew = main language, default RTL
\babelprovide[import]{english}        % (4) English secondary, for correct LTR runs
% Latin fallback registered BEFORE \babelfont so a lone English word in Hebrew
% prose does not tofu. Fonts resolved by FAMILY NAME (portable; no hard path).
\directlua{luaotfload.add_fallback("hebfb",{"DejaVuSerif:mode=node;"})}
% v3.22: Frank Ruhl Libre ships NO italic face. Instead of silently-upright
% \emph (v<=3.20) or bold-as-emphasis (v3.21), synthesize a slanted face with
% FakeSlant — real, VISIBLE italic emphasis with correct \emph nesting
% semantics (inner \emph toggles back upright). Verified empirically at
% 300dpi: slant clearly visible, 0 errors, 0 missing glyphs.
\babelfont{rm}[RawFeature={fallback=hebfb},
  ItalicFont={Frank Ruhl Libre}, ItalicFeatures={FakeSlant=0.22},
  BoldItalicFont={Frank Ruhl Libre Bold}, BoldItalicFeatures={FakeSlant=0.22}
  ]{Frank Ruhl Libre}
\babelfont{sf}[RawFeature={fallback=hebfb}]{Heebo}
\babelfont{tt}[RawFeature={fallback=hebfb}]{DejaVu Sans Mono}

\usepackage{amsmath,amssymb,mathtools}
\usepackage[margin=2.3cm,headheight=14pt]{geometry}
\linespread{1.3}                      % readable line spacing for study/annotation docs
\usepackage{xcolor}
\usepackage{mdframed}
\usepackage{listings}                 % code listings (assembly / C / pseudo-code)
\usepackage{tikz}
\usetikzlibrary{arrows.meta,positioning,calc,shapes.geometric,fit,backgrounds}
\usepackage{pgfplots}\pgfplotsset{compat=1.18}
\usepackage{quantikz}                 % quantum circuits (remove if unused)
\usepackage{booktabs}
\usepackage{longtable}
\usepackage{array}
\usepackage{float}                    % [H] = place table exactly here (reference doc: no drifting)
\usepackage{enumitem}
\usepackage{needspace}

% ============================================================================
% ★★ v3.4 — THE DIRECTION-ISOLATION TRIO (the heart of mixing Hebrew + Latin)
% Mental model:
%   \h{...}    Hebrew (RTL) island      — Hebrew text inside an LTR context
%                                          (TikZ nodes, the english `tikzpic`).
%   \en{...}   English (LTR) isolate    — Latin text / parentheses inside a
%                                          Hebrew (RTL) context. KEEPS PARENS
%                                          AND ORDER CORRECT.
%   \code{...} \en + monospace          — code inside Hebrew / inside a frametitle.
% CRITICAL: never wrap Latin in \h{} (it is for HEBREW). Inside a TikZ node,
% Latin in \h{} renders REVERSED (Datapath->htapataD, CPU->UPC) and parens
% mirror. Isolate every Latin/paren chunk with \en{} or \code{}. See
% references/figures-boxes-listings.md.
% ---- TEXT INSIDE MATH (v3.8) ----------------------------------------------
% Put ANY text that sits inside math in a \text{}. The ONLY thing that is
% invalid in math is a bare \h/\en (= \foreignlanguage) used DIRECTLY, outside
% \text -> "\rmfamily invalid in math mode". Inside \text:
%   Hebrew label           -> \text{...}            (bare Hebrew renders fine;
%                              \h is OPTIONAL/harmless: \text{\h{...}} also works
%                              and both render identically + scale in scripts)
%   plain Latin word        -> \text{...}            (no marker needed)
%   Latin WITH ( ) or [ ]   -> \text{\en{...}}       (\en forces LTR; a plain
%                              \text mirrors the parens in RTL -> "Spec (TNR(" )
%   f(Hebrew/text arg)      -> \text{f}(\text{...})  -- keep the parens at MATH
%                              level, NOT \text{f(arg)} (parens inside the box mirror)
%   code in math            -> \mathtt{...}  (or \texttt{}, math-guarded below)
% NEVER write \en{...} or \h{...} straight into $...$ / \[..\] (outside \text);
% always via \text{}. (\mbox{...} also works but does not scale in scripts.)
% ============================================================================
\newcommand{\h}[1]{\foreignlanguage{hebrew}{#1}}
\newcommand{\en}[1]{\foreignlanguage{english}{#1}}

% ============================================================================
% ★★ — fixes for defects that are INVISIBLE in the log and in the checks
% and only surface on a VISUAL RENDER of the PDF. (Each verified empirically.)
% ----------------------------------------------------------------------------
% v3.22: \emph is now REAL slanted emphasis (FakeSlant italic face defined on
% the \babelfont{rm} line above), replacing the v3.21 \emph->\textbf hack: it
% preserves \emph's toggle semantics and keeps bold free for its own role.
% Prefer bold-as-emphasis for your doc? one line: \renewcommand{\emph}[1]{\textbf{#1}}
% page number: force the digits LTR so they never reverse (18 -> 81) in the RTL
%   footer. That reversal is a NON-DETERMINISTIC bidi glitch (observed once in
%   ~80 pp); this LuaLaTeX build is not byte-reproducible, so the defensive
%   LTR wrap is the only reliable guard. Latin digits are always LTR anyway.
\renewcommand{\thepage}{\foreignlanguage{english}{\arabic{page}}}
% emergencystretch: Hebrew does not hyphenate, so an unbreakable LTR island
%   (\en{...} / an inline math run) landing at the margin can leave a small
%   (~2-13pt) overfull with no legal breakpoint. This rescues ONLY those
%   inline-prose tie-breaks; it is a ceiling used only on offending lines and
%   does NOT touch display-math / listing / table overflow (those still surface
%   in the log for the content-level overflow pass). Verified: baseline overfull
%   -> 0 with this set, with no change to pagination or underfull.
\setlength{\emergencystretch}{3em}
% em-dash: --- does NOT become — under Frank Ruhl (the font lacks tlig; verified
%   the width is identical with/without Ligatures=TeX, while the mechanism works
%   on Latin Modern). TYPE THE CHARACTER — (U+2014) DIRECTLY; likewise – (U+2013).
%   check_bidi_figures.py flags a --/--- that is adjacent to Hebrew.
% ============================================================================

% ---- STRUCTURAL DIRECTION-SAFETY for \texttt (v3.7) -------------------------
% A parenthesis, bracket, or $ INSIDE \texttt{} mirrors/migrates in RTL prose
% (MIPS 0($s0)->"(s0$)0", $rs->"rs$", arr[i]->"[arr[i") because \texttt sets the
% font but NOT the direction. Redefining \texttt to LTR-isolate its content fixes
% the WHOLE class at once (parens, brackets, $) so plain \texttt{lw $t0,0($s0)}
% is safe — no need to remember \code{}. Guarded by \ifmmode because
% \foreignlanguage is invalid in math (there, plain monospace is already LTR).
% NB: \texttt{Hebrew} stays tofu — the monospace font has no Hebrew anyway; put
% Hebrew in \h{}/normal text, never in \texttt{}. To disable, delete this block.
\let\origtexttt\texttt
\renewcommand{\texttt}[1]{\ifmmode\origtexttt{#1}\else\foreignlanguage{english}{\origtexttt{#1}}\fi}
\newcommand{\code}[1]{\ifmmode\origtexttt{#1}\else\foreignlanguage{english}{\origtexttt{#1}}\fi}

% ---- a thin decorative rule helper ----
\newcommand{\sep}{\par\vspace{2pt}\noindent\textcolor{cRule}{\rule{\linewidth}{0.6pt}}\par\vspace{4pt}}

% ---- colors ----
\definecolor{cBlue}{HTML}{1F4E79}   \definecolor{cBlueBg}{HTML}{EAF1F8}
\definecolor{cGreen}{HTML}{2E6B4F}  \definecolor{cGreenBg}{HTML}{EAF4EE}
\definecolor{cAmber}{HTML}{9C6A00}  \definecolor{cAmberBg}{HTML}{FBF3E2}
\definecolor{cPurple}{HTML}{5B2C6F} \definecolor{cPurpleBg}{HTML}{F3ECF7}
\definecolor{cTeal}{HTML}{0E6E6E}   \definecolor{cTealBg}{HTML}{E6F4F4}
\definecolor{cRed}{HTML}{8E2B2B}    \definecolor{cRedBg}{HTML}{FBEAEA}
\definecolor{cGray}{HTML}{555555}   \definecolor{cRule}{HTML}{1F4E79}

% ---- math operators ----
\DeclareMathOperator{\Tr}{Tr}
\DeclareMathOperator{\rank}{rank}
\DeclareMathOperator{\diag}{diag}
\DeclareMathOperator{\sgn}{sgn}

% ---- bra-ket (quantikz defines \ket/\bra/\braket; add the rest safely) ----
\providecommand{\ket}[1]{\left|#1\right\rangle}
\providecommand{\bra}[1]{\left\langle#1\right|}
\providecommand{\braket}[2]{\left\langle#1\middle|#2\right\rangle}
\newcommand{\dyad}[1]{\left|#1\right\rangle\!\left\langle#1\right|}
% v3.21: general outer product |a><b| and expectation value <A> — agents reach
% for these by reflex (esp. \expval); quantikz defines neither, so their absence
% threw "Undefined control sequence" in a real build. \providecommand = safe.
\providecommand{\ketbra}[2]{\left|#1\right\rangle\!\left\langle#2\right|}
\providecommand{\expval}[1]{\left\langle#1\right\rangle}
\newcommand{\abs}[1]{\left|#1\right|}
\newcommand{\norm}[1]{\left\lVert#1\right\rVert}
\newcommand{\half}{\tfrac12}
\newcommand{\R}{\mathbb{R}}\newcommand{\C}{\mathbb{C}}\newcommand{\Z}{\mathbb{Z}}
% identity operator. NB: \mathbb{1} silently renders a WRONG glyph (msbm has no
% blackboard digits). Use \mathbb{I} = 𝕀 (works with amssymb, no extra package).
% Do NOT "fix" this back to \mathbb{1}.
\newcommand{\id}{\mathbb{I}}

% ============================================================================
% Colored boxes — mdframed (NOT tcolorbox, which breaks under bidi).
% ★★ v3.4: innerleftmargin/innerrightmargin give code listings room so their
% frame/line-numbers do not spill out the box border (see listing styles below).
% ============================================================================
\newmdenv[linecolor=cBlue,backgroundcolor=cBlueBg,linewidth=1.1pt,
  frametitlebackgroundcolor=cBlue,frametitlefontcolor=white,
  innertopmargin=6pt,innerbottommargin=7pt,skipabove=9pt,skipbelow=9pt,
  innerleftmargin=9pt,innerrightmargin=9pt,roundcorner=3pt]{defbox}   % הגדרה
\newmdenv[linecolor=cGreen,backgroundcolor=cGreenBg,linewidth=1.1pt,
  frametitlebackgroundcolor=cGreen,frametitlefontcolor=white,
  innertopmargin=6pt,innerbottommargin=7pt,skipabove=9pt,skipbelow=9pt,
  innerleftmargin=9pt,innerrightmargin=9pt,roundcorner=3pt]{thmbox}   % משפט/טענה
\newmdenv[linecolor=cAmber,backgroundcolor=cAmberBg,linewidth=1.1pt,
  frametitlebackgroundcolor=cAmber,frametitlefontcolor=white,
  innertopmargin=6pt,innerbottommargin=7pt,skipabove=9pt,skipbelow=9pt,
  innerleftmargin=9pt,innerrightmargin=9pt,roundcorner=3pt]{notebox}  % הערה
\newmdenv[linecolor=cPurple,backgroundcolor=cPurpleBg,linewidth=1.1pt,
  frametitlebackgroundcolor=cPurple,frametitlefontcolor=white,
  innertopmargin=6pt,innerbottommargin=7pt,skipabove=9pt,skipbelow=9pt,
  innerleftmargin=9pt,innerrightmargin=9pt,roundcorner=3pt]{exbox}    % דוגמה
\newmdenv[linecolor=cRed,backgroundcolor=cRedBg,linewidth=1.1pt,
  frametitlebackgroundcolor=cRed,frametitlefontcolor=white,
  innertopmargin=6pt,innerbottommargin=7pt,skipabove=9pt,skipbelow=9pt,
  innerleftmargin=9pt,innerrightmargin=9pt,roundcorner=3pt]{warnbox}  % טעות נפוצה
\newmdenv[linecolor=cTeal,backgroundcolor=cTealBg,linewidth=1.1pt,
  frametitlebackgroundcolor=cTeal,frametitlefontcolor=white,
  innertopmargin=6pt,innerbottommargin=7pt,skipabove=9pt,skipbelow=9pt,
  innerleftmargin=9pt,innerrightmargin=9pt,roundcorner=3pt]{keybox}   % רעיון מרכזי

% ============================================================================
% Unified TABLE presentation: subtle uniform slate frame + auto-numbered
% caption rendered INSIDE the frame. Use:
%   \begin{tablebox}\centering\small  <tabular>  \tabcap{<description>}\end{tablebox}
% Replaces the buggy \caption* (which emitted a stray "טבלה N : *" line) and
% gives every table one identical frame + numbered caption "טבלה N.M.".
% Optional intro text may be placed at the TOP, before the tabular, inside the box.
% ============================================================================
\definecolor{cTblLine}{HTML}{B9C4D2}  \definecolor{cTblBg}{HTML}{FAFBFC}
% v3.23: class-aware — under `article` there is no chapter counter (the promised
% report->article switch used to die with "No counter 'chapter' defined").
\makeatletter
\@ifundefined{c@chapter}{%
  \newcounter{tbl}[section]\renewcommand{\thetbl}{\arabic{tbl}}%
}{%
  \newcounter{tbl}[chapter]\renewcommand{\thetbl}{\thechapter.\arabic{tbl}}%
}
\makeatother
\newmdenv[linecolor=cTblLine,backgroundcolor=cTblBg,linewidth=0.9pt,
  innertopmargin=9pt,innerbottommargin=8pt,skipabove=11pt,skipbelow=11pt,
  innerleftmargin=11pt,innerrightmargin=11pt,roundcorner=4pt]{tablebox}
\newcommand{\tabcap}[1]{\refstepcounter{tbl}\par\addvspace{6pt}%
  {\centering\small\textcolor{cGray}{\textbf{טבלה~\thetbl.}}\enspace #1\par}}

% ============================================================================
% LTR islands.  EVERY tikzpicture / quantikz MUST live inside one, or the whole
% figure's coordinates and Latin labels flip. (Hebrew node text still uses \h{}.)
% ============================================================================
\newenvironment{qcirc}{\par\medskip\begin{otherlanguage}{english}\begin{center}}
  {\end{center}\end{otherlanguage}\par\medskip}
\newenvironment{tikzpic}{\par\medskip\begin{otherlanguage}{english}\begin{center}}
  {\end{center}\end{otherlanguage}\par\medskip}

% ============================================================================
% Code listings.  ★★ v3.4 THE BOX-OVERFLOW FIX:
%   resetmargins=true   → margins measured from the ENCLOSING box, not the page
%                         (without it, a listing inside an mdframed box spills
%                          its frame + line-numbers out past the box border).
%   framexleftmargin=0  → the frame (left bar) does not extend left of the code.
%   xleftmargin keeps a small indent so the line numbers have room INSIDE.
% asmblock wraps the listing in an LTR island (code is LTR).
% ============================================================================
\lstdefinestyle{asm}{
  basicstyle=\ttfamily\small,
  backgroundcolor=\color{cBlueBg!55},
  frame=leftline,framerule=2.5pt,rulecolor=\color{cBlue},
  numbers=left,numberstyle=\tiny\color{cGray},numbersep=8pt,
  xleftmargin=2.2em,framexleftmargin=0pt,framexrightmargin=0pt,resetmargins=true,
  breaklines=true,columns=fullflexible,keepspaces=true,
  showstringspaces=false,aboveskip=8pt,belowskip=8pt,
  commentstyle=\color{cGreen},
}
\lstdefinestyle{cstyle}{
  language=C,basicstyle=\ttfamily\small,
  backgroundcolor=\color{cGreenBg!55},
  frame=leftline,framerule=2.5pt,rulecolor=\color{cGreen},
  numbers=left,numberstyle=\tiny\color{cGray},numbersep=8pt,
  xleftmargin=2.2em,framexleftmargin=0pt,framexrightmargin=0pt,resetmargins=true,
  breaklines=true,columns=fullflexible,keepspaces=true,
  showstringspaces=false,keywordstyle=\color{cBlue}\bfseries,
  commentstyle=\color{cGray}\itshape,aboveskip=8pt,belowskip=8pt,
}
\newenvironment{asmblock}{\par\begin{otherlanguage}{english}}{\end{otherlanguage}\par}

\usepackage[hidelinks,unicode]{hyperref}   % LAST, after colors. unicode = Hebrew bookmarks.

\setcounter{tocdepth}{1}
\setcounter{secnumdepth}{2}

\usepackage{graphicx}

\title{מבנה המחשב הקוונטי}
\author{ראווה ל־hebrew-lualatex-pdf — לקט צפוף}
\date{}

\begin{document}
\begin{titlepage}
\centering\vspace*{4.5cm}
{\LARGE\bfseries מבנה המחשב הקוונטי\par}
\vspace{1.2cm}
{\large לקט צפוף — ארכיטקטורה קלאסית, מעגלים קוונטיים, והחיבור ביניהם\par}
\vspace{0.6cm}
{\large ראווה ל־\en{hebrew-lualatex-pdf}\par}
\vfill
\end{titlepage}

\tableofcontents
\clearpage

% ═══════════════════════════════════════════════════════════════
\part{צנרת ומסופים}
% ═══════════════════════════════════════════════════════════════

\chapter{חמשת שלבי הצנרת}
\label{ch:pipeline}

מעבד מודרני מפרק הוראה לחמישה שלבים עוקבים. כל מחזור שעון מקדם הוראה אחת לכל שלב,
כך שבמצב יציב יוצאת הוראה מושלמת בכל מחזור.\footnote{זהו ה־\en{throughput}; ה־\en{latency}
של הוראה בודדת נשאר חמישה מחזורים.}

\begin{defbox}[frametitle={הגדרה — צנרת מודולרית}]
צנרת \en{pipeline} היא פיצול ביצוע הוראה לרצף שלבים חופפים בזמן.
חמשת השלבים הקלאסיים: \en{IF, ID, EX, MEM, WB}.
\end{defbox}

\begin{equation}
\label{eq:cpi-ideal}
\mathrm{CPI}_{\mathrm{ideal}} = 1
\end{equation}
משוואה~\eqref{eq:cpi-ideal} היא האידיאל. בפועל, סכסוכים מעלים את \en{CPI} מעל אחד.

\begin{tikzpic}
\begin{tikzpicture}[font=\small,node distance=6mm,>={Latex[length=2.2mm]},
  box/.style={draw,thick,rounded corners,align=center,minimum width=1.9cm,minimum height=1.15cm,fill=cBlueBg}]
  \node[box] (if)  {\en{IF}\\שליפה};
  \node[box,right=of if]  (id)  {\en{ID}\\פענוח};
  \node[box,right=of id]  (ex)  {\en{EX}\\ביצוע};
  \node[box,right=of ex]  (mem) {\en{MEM}\\זיכרון};
  \node[box,right=of mem,fill=cGreenBg] (wb) {\en{WB}\\כתיבה};
  \draw[->,thick] (if) -- (id);
  \draw[->,thick] (id) -- (ex);
  \draw[->,thick] (ex) -- (mem);
  \draw[->,thick] (mem) -- (wb);
\end{tikzpicture}
\end{tikzpic}
\begin{center}\small איור — חמשת שלבי הצנרת הקלאסית\end{center}

\begin{tablebox}\centering\small
\begin{tabular}{@{}lll@{}}
\toprule
\textbf{שלב} & \textbf{שם} & \textbf{פעולה מרכזית}\\
\midrule
\en{IF}  & שליפת הוראה     & קריאה מ־\en{IMEM} לפי \en{PC}\\
\en{ID}  & פענוח           & קריאת אוגרים, הרחבת מיידי\\
\en{EX}  & ביצוע           & חישוב ב־\en{ALU} / חישוב כתובת\\
\en{MEM} & גישה לזיכרון    & \code{lw/sw} ל־\en{DMEM}\\
\en{WB}  & כתיבה לאחור     & עדכון קובץ האוגרים\\
\bottomrule
\end{tabular}
\tabcap{חמשת השלבים ותפקידם.}
\end{tablebox}

\chapter{סכסוכי נתונים וקידום}
\label{ch:hazards}

סכסוך נתונים נוצר כשהוראה זקוקה לתוצאה של קודמתה שעדיין לא הגיעה ל־\en{WB}.
הפתרון הנפוץ: \en{forwarding} ממעקפי \en{EX/MEM} ו־\en{MEM/WB}.

\begin{keybox}[frametitle={רעיון מפתח — קידום}]
במקום להמתין ל־\en{WB}, מעבירים את התוצאה ישירות מיציאת ה־\en{ALU}
(או מ־\en{MEM}) לקלט של ההוראה התלויה.
\end{keybox}

\begin{warnbox}[frametitle={מלכודת — \code{load-use}}]
גם עם קידום מלא, \code{lw} ואחריו שימוש מיידי דורשים \en{stall} של מחזור אחד —
הנתונים מגיעים רק בסוף \en{MEM}.
\end{warnbox}

\begin{tikzpic}
\begin{tikzpicture}[font=\small,>={Latex[length=2.2mm]},
  n/.style={draw,thick,rounded corners,align=center,minimum width=3.2cm,minimum height=1.05cm}]
  \node[n,fill=cBlueBg] (q) {תלות בנתונים?};
  \node[n,fill=cGreenBg,below left=1.4cm and 0.6cm of q] (fwd) {קידום מספיק};
  \node[n,fill=cRedBg,below right=1.4cm and 0.6cm of q] (stall) {\en{load-use stall}};
  \draw[->,thick] (q) -- node[left,font=\footnotesize] {לא־טעינה} (fwd);
  \draw[->,thick] (q) -- node[right,font=\footnotesize] {\en{lw $\to$ use}} (stall);
\end{tikzpicture}
\end{tikzpic}
\begin{center}\small איור — עץ החלטה לטיפול בתלות\end{center}

\begin{exbox}[frametitle={דוגמה — רצף עם תלות}]
הרצף הבא יוצר תלות \en{RAW} על \code{\$t0}:
\begin{asmblock}
\begin{lstlisting}[style=asm]
add  $t0, $s0, $s1    # t0 = s0 + s1
sub  $t2, $t0, $s2    # needs t0 — forward from EX/MEM
lw   $t3, 0($t0)      # address from t0 — also forward
addi $t4, $t3, 4      # load-use: stall required
\end{lstlisting}
\end{asmblock}
\end{exbox}

\begin{equation}
\label{eq:cpu-time}
T = \mathrm{IC}\times\mathrm{CPI}\times T_{\mathrm{clk}}
\end{equation}
זמן הריצה~\eqref{eq:cpu-time} הוא המכפלה הקלאסית. שיפור צנרת מקטין בעיקר את \en{CPI}
ואת $T_{\mathrm{clk}}$.

% ═══════════════════════════════════════════════════════════════
\part{היררכיית זיכרון}
% ═══════════════════════════════════════════════════════════════

\chapter{מטמון וממוצע גישה}
\label{ch:cache}

היררכיית הזיכרון מנצלת מקומיות. זמן הגישה האפקטיבי למערכת עם מטמון ברמה אחת:

\begin{equation}
\label{eq:tavg}
T_{\mathrm{avg}} = T_{\mathrm{hit}} + r_{\mathrm{miss}}\cdot T_{\mathrm{miss}}
\end{equation}
במשוואה~\eqref{eq:tavg}: \en{Hit rate $=1-r_{\mathrm{miss}}$}.
המשמעות: גם \en{miss rate} קטן מכביד אם
$\underbrace{T_{\mathrm{miss}}}_{\text{זמן החמצה}}$ גדול.

\begin{keybox}[frametitle={נוסחת העוגן — \en{AMAT}}]
זמן הגישה הממוצע $\overbrace{T_{\mathrm{avg}}}^{\text{\en{AMAT}}}$ הוא המדד המרכזי
לתכנון מטמון. כל שיפור ב־\en{hit time} או ב־\en{miss rate} נמדד בו.
\end{keybox}

\begin{tablebox}\centering\small
\begin{tabular}{@{}llll@{}}
\toprule
\textbf{\en{Hit}} & \textbf{\en{Miss rate}} & \textbf{\en{Miss penalty}} & \textbf{\en{AMAT}}\\
\midrule
$1$ מחזור   & $5\%$ & $100$ & $6.0$\\
$1$ מחזור   & $2\%$ & $100$ & $3.0$\\
$2$ מחזורים & $1\%$ & $100$ & $3.0$\\
\bottomrule
\end{tabular}
\tabcap{דוגמאות מספריות ל־\en{AMAT}.}
\end{tablebox}

\begin{tikzpic}
\begin{tikzpicture}
\begin{axis}[
  width=11cm, height=6.2cm,
  xlabel={גודל מטמון (KB)},
  ylabel={שיעור החמצה (\%)},
  xmin=0, xmax=34, ymin=0, ymax=20,
  legend style={at={(0.97,0.97)},anchor=north east,font=\small},
  grid=major, grid style={gray!20},
]
\addplot[thick,mark=*] coordinates {(1,18)(2,14)(4,10)(8,7)(16,5)(32,4)};
\addlegendentry{ישיר}
\addplot[thick,dashed,mark=square*] coordinates {(1,14)(2,10)(4,7)(8,4.5)(16,3)(32,2.2)};
\addlegendentry{סט־4}
\end{axis}
\end{tikzpicture}
\end{tikzpic}
\begin{center}\small איור — שיעור החמצה כפונקציה של גודל ואסוציאטיביות\end{center}

\begin{notebox}[frametitle={הערה — שלוש ה־C}]
מסווגים החמצות ל־\en{Compulsory, Capacity, Conflict}.
הגדלת בלוק מקטינה compulsory; הגדלת מטמון מקטינה capacity; אסוציאטיביות מקטינה conflict.
\end{notebox}

\chapter{טבלת סמלים לזיכרון}
\label{ch:cache-glossary}

\begin{tablebox}\centering\small
\begin{tabular}{@{}rp{9.5cm}@{}}
\toprule
\textbf{סמל} & \textbf{משמעות}\\
\midrule
$T_{\mathrm{hit}}$  & זמן גישה בפגיעה\\
$r_{\mathrm{miss}}$ & שיעור החמצות\\
$T_{\mathrm{miss}}$ & עונש החמצה (מחזורים)\\
\en{AMAT}           & זמן גישה ממוצע\\
\en{LRU}            & מדיניות החלפה: הכי פחות בשימוש לאחרונה\\
\bottomrule
\end{tabular}
\tabcap{מילון סמלים — מטמון.}
\end{tablebox}

% ═══════════════════════════════════════════════════════════════
\part{חישוב קוונטי}
% ═══════════════════════════════════════════════════════════════

\chapter{הקיוביט ומרחב המצבים}
\label{ch:qubit}

קיוביט הוא וקטור יחידה במרחב הילברט דו־ממדי. מצב כללי הוא סופרפוזיציה של בסיס החישוב,
עם משרעות מרוכבות.\footnote{ההסתברות היא ריבוע הערך המוחלט של המשרעת — לא המשרעת עצמה.}

\begin{equation}
\label{eq:state}
\ket\psi = \alpha\ket0 + \beta\ket1,\qquad |\alpha|^2+|\beta|^2=1
\end{equation}
משוואה~\eqref{eq:state} מגדירה מצב טהור. על כדור בלוך:
\[
\ket\psi = \cos\tfrac\theta2\ket0 + e^{i\varphi}\sin\tfrac\theta2\ket1.
\]

\begin{defbox}[frametitle={הגדרה — כדור בלוך}]
כל מצב קיוביט טהור מתאים לנקודה על ספירת היחידה.
הקוטב הצפוני הוא $\ket0$, הדרומי $\ket1$.
\end{defbox}

\begin{tablebox}\centering\small
\begin{tabular}{@{}ccc@{}}
\toprule
\textbf{שער} & \textbf{מטריצה} & \textbf{פעולה}\\
\midrule
$X$ & $\begin{pmatrix}0&1\\1&0\end{pmatrix}$ & היפוך סיבית\\
$Z$ & $\begin{pmatrix}1&0\\0&-1\end{pmatrix}$ & היפוך פאזה\\
$H$ & $\tfrac1{\sqrt2}\begin{pmatrix}1&1\\1&-1\end{pmatrix}$ & סופרפוזיציה\\
$S$ & $\begin{pmatrix}1&0\\0&i\end{pmatrix}$ & פאזה $\pi/2$\\
\bottomrule
\end{tabular}
\tabcap{שערים בסיסיים של קיוביט בודד.}
\end{tablebox}

\begin{equation}
\label{eq:hgate}
H\ket0 = \tfrac1{\sqrt2}(\ket0+\ket1) = \ket{+}
\end{equation}
שער הדמרד~\eqref{eq:hgate} יוצר סופרפוזיציה שווה:
$\underbrace{H\ket0}_{\text{סופרפוזיציה שווה}}$.

\chapter{שזירה ומצבי בל}
\label{ch:bell}

מצב בל הוא המצב השזור המקסימלי הפשוט ביותר לשני קיוביטים:
\begin{equation}
\label{eq:bell}
\ket{\Phi^+} = \tfrac1{\sqrt2}(\ket{00}+\ket{11})
\end{equation}

\begin{qcirc}
\begin{quantikz}
\lstick{$\ket0$} & \gate{H} & \ctrl{1} & \meter{} \\
\lstick{$\ket0$} & \qw      & \targ{}  & \meter{}
\end{quantikz}
\end{qcirc}
\begin{center}\small איור — הכנת מצב בל ומדידה\end{center}

\begin{thmbox}[frametitle={משפט — אי־שכפול}]
לא קיים אופרטור אוניטרי $U$ המעתיק מצב שרירותי:
$U(\ket\psi\otimes\ket0)=\ket\psi\otimes\ket\psi$ לא ייתכן לכל $\ket\psi$.
\end{thmbox}

\noindent\textbf{הוכחה (בקצרה).}

נניח שקיים $U$ כזה לשני מצבים. אז $\braket{\psi}{\phi}=\braket{\psi}{\phi}^2$,
ולכן $\braket{\psi}{\phi}\in\{0,1\}$ — סתירה לליניאריות על מצב שרירותי. $\qquad\blacksquare$

\chapter{טלפורטציה קוונטית}
\label{ch:teleport}

טלפורטציה מעבירה מצב לא ידוע באמצעות שזירה משותפת ושתי סיביות קלאסיות —
בלי לשלוח את הקיוביט עצמו.

\begin{keybox}[frametitle={פרוטוקול הטלפורטציה}]
\textbf{(1)} אליס ובוב חולקים $\ket{\Phi^+}$ ממשוואה~\eqref{eq:bell}.\\
\textbf{(2)} אליס מודדת בבסיס בל את הקיוביט שלה ואת $\ket\psi$, ושולחת $2$ סיביות.\\
\textbf{(3)} בוב מפעיל $X$ ו/או $Z$ לפי התוצאה ומשחזר את $\ket\psi$.
\end{keybox}

\begin{qcirc}
\resizebox{0.92\textwidth}{!}{%
\begin{quantikz}
\lstick{$\ket\psi$} & \qw & \qw & \ctrl{1} & \gate{H} & \meter{} \\
\lstick{$\ket0$} & \gate{H} & \ctrl{1} & \targ{} & \qw & \meter{} \\
\lstick{$\ket0$} & \qw & \targ{} & \qw & \qw & \gate{X}\gate{Z}
\end{quantikz}%
}
\end{qcirc}
\begin{center}\small איור — מעגל טלפורטציה (פישוט סכמטי)\end{center}

\begin{keybox}[frametitle={לקח}]
שזירה + ערוץ קלאסי = העברת מצב קוונטי.
אי־שכפול מבטיח שאי אפשר להעתיק בלי להרוס את המקור.
\end{keybox}

% ═══════════════════════════════════════════════════════════════
\part{החיבור: בקרה קלאסית על מעגל קוונטי}
% ═══════════════════════════════════════════════════════════════

\chapter{בקרת משוב היברידית}
\label{ch:hybrid}

מערכת קוונטית מעשית נשלטת ע״י מחשב קלאסי: תזמון פולסים, קריאת מדידות,
ועדכון פרמטרים באלגוריתמים וריאציוניים (\en{VQE, QAOA}).

\begin{tikzpic}
\begin{tikzpicture}[font=\small,node distance=1.6cm,>={Latex[length=2.4mm]},
  box/.style={draw,thick,rounded corners,align=center,minimum width=2.8cm,minimum height=1.2cm,fill=cBlueBg}]
  \node[box] (cpu)  {מעבד \en{(CPU)}};
  \node[box,right=of cpu] (ctrl) {בקר \en{AWG}};
  \node[box,right=of ctrl,fill=cTealBg] (qpu) {שבב \en{(QPU)}};
  \draw[<->,thick] (cpu) -- node[above,font=\footnotesize] {פרמטרים} (ctrl);
  \draw[->,thick] (ctrl) -- node[above,font=\footnotesize] {פולסים} (qpu);
  \draw[<-,thick,dashed] (cpu.south) to[bend right=20] node[below,font=\footnotesize] {מדידות} (qpu.south);
\end{tikzpicture}
\end{tikzpic}
\begin{center}\small איור — לולאת בקרה היברידית\end{center}

\begin{notebox}[frametitle={מוסכמה}]
בקוד הבקרה משתמשים בהערות באנגלית בלבד. עברית בפרוזה; מזהים טכניים ב־\code{code}.
\end{notebox}

\begin{exbox}[frametitle={דוגמה — שליפת תוצאת מדידה}]
\begin{asmblock}
\begin{lstlisting}[style=cstyle]
# Classical side: collect bitstrings from QPU shots
counts = backend.run(circuit, shots=1024).result().get_counts()
p00 = counts.get("00", 0) / 1024
# Estimate <Z x Z> from parity of bitstrings
expval = sum(((-1)**bin(int(b,2)).count("1")) * c
             for b, c in counts.items()) / 1024
\end{lstlisting}
\end{asmblock}
\end{exbox}

\begin{tablebox}\centering\small
\begin{otherlanguage}{english}
\begin{tabular}{@{}lll@{}}
\toprule
\textbf{bitstring} & $\boldsymbol{Z\otimes Z}$ & \h{תרומה ל־}$\langle Z\otimes Z\rangle$\\
\midrule
\texttt{00} & $(+1)$ & \h{חיובית}\\
\texttt{01} & $(-1)$ & \h{שלילית}\\
\texttt{10} & $(-1)$ & \h{שלילית}\\
\texttt{11} & $(+1)$ & \h{חיובית}\\
\bottomrule
\end{tabular}
\end{otherlanguage}
\tabcap{טבלת אמת — פריטי מדידה בשני קיוביטים.}
\end{tablebox}

\chapter{התמונה הגדולה}
\label{ch:bigpic}

שלושה עקרונות מקשרים את החלקים:
\begin{enumerate}
\item[(\textbf{א})] צנרת קלאסית ממקסמת \en{throughput} תחת תלויות —
  קידום ו־\en{stall} הם מחיר התלויות.
\item[(\textbf{ב})] מטמון ממיר מקומיות ל־\en{AMAT} נמוך —
  הנוסחה~\eqref{eq:tavg} היא מצפן התכנון.
\item[(\textbf{ג})] מעגל קוונטי ממיר שזירה ליכולת שלא קיימת קלאסית —
  טלפורטציה ואי־שכפול הם הגבולות.
\end{enumerate}

\begin{keybox}[frametitle={סיכום — נוסחת הזמן הקלאסית עדיין שולטת בבקרה}]
גם במערכת היברידית, זמן הריצה של לולאת הבקרה נמדד ב־\eqref{eq:cpu-time}.
ה־\en{QPU} הוא מאיץ; הצוואר הוא לעתים קרובות ה־\en{CPU} והערוץ.
\end{keybox}

% ═══════════════════════════════════════════════════════════════
\part{תרגילים פתורים}
% ═══════════════════════════════════════════════════════════════

\chapter{תרגילים}
\label{ch:exercises}

\begin{exbox}[frametitle={תרגיל 1 — חישוב \en{AMAT}}]
\textbf{(א)} מטמון עם \en{hit time $=1$, miss rate $=0.04$, miss penalty $=80$}.
חשב \en{AMAT}.\\[4pt]
\textbf{(ב)} אם מגדילים את המטמון ו־\en{miss rate} יורד ל־$0.02$ אבל
\en{hit time} עולה ל־$2$, האם כדאי?
\end{exbox}

\begin{keybox}[frametitle={פתרון 1},nobreak=true]
\textbf{פתרון (א).}\\
לפי~\eqref{eq:tavg}:
\[ T_{\mathrm{avg}}=1+0.04\cdot 80=4.2 \]
מחזורים.\\[4pt]
\textbf{פתרון (ב).}\\
\[ T_{\mathrm{avg}}=2+0.02\cdot 80=3.6 \]
כן, כדאי: ירידה מ־$4.2$ ל־$3.6$.\\[4pt]
\textbf{לקח.} הקטנת \en{miss rate} יכולה לפצות על עלייה ב־\en{hit time} —
בודקים תמיד ב־\en{AMAT}.
\end{keybox}

\begin{exbox}[frametitle={תרגיל 2 — הסתברות במצב בל}]
נתון $\ket{\Phi^+}$ ממשוואה~\eqref{eq:bell}. מהי $P(00)$? מהי $P(01)$?
\end{exbox}

\begin{keybox}[frametitle={פתרון 2},nobreak=true]
\textbf{פתרון.}
\[ P(00)=\bigl|\braket{00}{\Phi^+}\bigr|^2=\tfrac12 \]
$P(01)=0$ — אין רכיב $\ket{01}$ במצב.\\[4pt]
\textbf{לקח.} במצב בל אידיאלי המדידות מתואמות לחלוטין; תוצאות מנוגדות מתאפסות.
\end{keybox}

\begin{exbox}[frametitle={תרגיל 3 — סכסוך \code{load-use}}]
עבור הרצף:
\begin{asmblock}
\begin{lstlisting}[style=asm]
lw   $t0, 0($s0)
add  $t1, $t0, $s1
\end{lstlisting}
\end{asmblock}
כמה מחזורי \en{stall} נדרשים עם קידום מלא? בלי קידום?
\end{exbox}

\begin{keybox}[frametitle={פתרון 3},nobreak=true]
\textbf{פתרון.}\\
עם קידום מלא: מחזור \en{stall} אחד (\en{load-use}).\\
בלי קידום: שני מחזורים (המתנה עד אחרי \en{WB}).\\[4pt]
\textbf{לקח.} קידום לא מבטל את כל העצירות — רק את רובן.
\end{keybox}

\chapter{נספח — זהויות שימושיות}
\label{ch:appendix}

זהויות בסיסיות לשימוש חוזר:
\[ HXH = Z,\qquad HZH = X,\qquad H^2 = I \]
\[
\ket{\Phi^\pm}=\tfrac1{\sqrt2}(\ket{00}\pm\ket{11}),\qquad
\ket{\Psi^\pm}=\tfrac1{\sqrt2}(\ket{01}\pm\ket{10})
\]

\noindent ציטוט סיום: המסמך הזה הוא ״מבחן קצה״ לצינור —
מתמטיקה, דיאגרמות עם לטינית מעורבת, מעגלים, גרפים, קופסאות, טבלאות וקוד —
בלי tofu ובלי היפוכי כיוון.

\end{document}
